%
% komplexer-summand.tex -- 
%
% (c) 2021 Prof Dr Andreas Müller, OST Ostschweizer Fachhochschule
%
\documentclass[tikz]{standalone}
\usepackage{amsmath}
\usepackage{times}
\usepackage{txfonts}
\usepackage{pgfplots}
\usepackage{csvsimple}
\usetikzlibrary{arrows,intersections,math,calc}
\definecolor{darkred}{rgb}{0.8,0,0}
\definecolor{darkgreen}{rgb}{0,0.6,0}
\begin{document}
\def\skala{1.05}
\begin{tikzpicture}[>=latex,thick,scale=\skala]

	\draw[->] (-0.4,0) -- (4.7,0) coordinate[label={$\operatorname{Re}$}];
	\draw[->] (0,-2.5) -- (0,2.7) coordinate[label={left:$\operatorname{Im}$}];
	\draw[dashed] (0,0) circle[radius=2.0];
	\draw[->,very thick] (0,0) -- (1.55,-1.26)
		node[midway,below left] {$n^{-\sigma}$};
	\draw (0.75,0) arc[start angle=0,end angle=-39,radius=0.75];
	\node at (1.35,-0.36) {$-t\log n$};
	\fill (1.55,-1.26) circle[radius=1.5pt];
	\node[right] at (1.55,-1.26)
		{$n^{-\sigma}e^{-it\log n}$};

\end{tikzpicture}
\end{document}

